\scanpage{116}

soit ainsi décrit par sa structure de graphe déjà -- d'ailleurs même sans supposer la structure orientable ni régulière, ce n'est pas évident qu'il y ait un contre-exemple. \deleted{Mais si --} Mais si -- dans l'\uline{icosaèdre gauche}, (régulier, mais pas orientable), on  $S$ de cardinal 6, \deleted{\unsure{tte}} \deleted{\unsure{tout|toute}} $a\in\unsure{\mathbb P_2|\mathcal P_2|\mathfrak P_2}(S)$ est une arête (il y en a 15), mais seulement la moitié des 20 éléments $f\in\mathbb P_3(S)$ sont des faces\ldots.

% slanted
\margin{\uline{Digression sur relation entre $A_0$ et $F_0$}}

On peut essayer de voir, pour une structure de graphe \inserted{strict}, $\Gamma_0=(S_{\unsure{0|2}},A_0)$, à quelles conditions l'ens $F_0\subset\mathbb P_3(S_{\unsure{0}})$ qu'il \deleted{\unsure{en}} définit est une structure de 2-polyèdre. On trouve \deleted{\unsure{ce|ceci}}

\[
\begin{aligned}
&\deleted{\unsure{(40)}}\\
&A1)\quad \forall\,\underset{\substack{\shortparallel\\\{s,t\}}}{a}\in A_0,\ 
 \text{il existe exactement deux éléments}\\
&\qquad u\in S,\ \text{tels que l'on ait}\\
&\qquad \{s,u\},\{t,u\}\in A_0\quad \text{i.e. adjacents}\\
&\qquad \text{à la fois aux deux extrémités}\\
&\qquad \text{de l'arête }a\quad\text{i.e. }A_0(s)\cap A_0(t)\text{ est}\\
&\qquad \text{de cardinal }2.
\end{aligned}\tag*{$(39)$}
\]

\margin{\uline{NB} le fait qu'il existe au, moins \uline{un} tel $u$ indique déjà que $A_0$ se déduit de $F_0$ par $A_0=\{a\in\mathbb P_2(S)\mid\exists f\in F_{\unsure{3|0}}\ f\supset a\}$}

\[
\begin{aligned}
&A2)\quad \forall s\in S,\quad\text{la structure de contour dont l'ens.}\\
&\qquad\text{des sommets est }\inserted{A_0(s)=}\operatorname{Et}_{\Gamma_0}(s)\quad\text{(ens.}\\
&\qquad\text{des sommets adjacents), et où}\\
&\qquad\text{les arêtes sont les } a\in\mathbb P_2(A_0(s))\\
&\qquad\text{telles que }a\in A_0,\quad\text{est \uline{connexe}}\\
&\qquad\text{non vide \deleted{(\uline{NB} si elle est non v.}}
\end{aligned}
\]

\margin{\uline{NB} Si $A_0(s)\ne\varnothing$, il résulte déjà de $A1$ que $\operatorname{Card}A_0(s)\geqslant3$}
% archive sheet number: 115
